%%======================================================================
%%  MT-1 REDO: Literature Extraction & Convention Fixing (v2)
%%        for Black Hole Entropy from the Spectral Action
%%
%%  This document collects ALL formulas needed to derive the BH entropy
%%  from the SCT spectral action, compute c_log, and verify the four
%%  laws of BH thermodynamics.  Corrected from v1 (sign error in the
%%  Sen formula caught by L-R audit).
%%
%%  RESEARCH DOCUMENT: formulas extracted from cited papers only.
%%  No original derivations --- gaps noted explicitly.
%%======================================================================
\documentclass[12pt,a4paper]{article}

%% ---- packages -------------------------------------------------------
\usepackage[utf8]{inputenc}
\usepackage[T1]{fontenc}
\usepackage{lmodern}
\usepackage{amsmath,amssymb,amsthm,mathtools}
\usepackage[margin=2.5cm]{geometry}
\usepackage{booktabs}
\usepackage{enumitem}
\usepackage{hyperref}
\usepackage{xcolor}
\usepackage{fancyhdr}
\usepackage{longtable}
\usepackage{tcolorbox}

%% ---- theorem-like environments --------------------------------------
\newtheorem{proposition}{Proposition}[section]
\newtheorem{lemma}[proposition]{Lemma}
\newtheorem{corollary}[proposition]{Corollary}
\newtheorem{theorem}[proposition]{Theorem}
\theoremstyle{definition}
\newtheorem{definition}[proposition]{Definition}
\newtheorem{convention}[proposition]{Convention}
\theoremstyle{remark}
\newtheorem{remark}[proposition]{Remark}

%% ---- custom commands ------------------------------------------------
\newcommand{\Tr}{\operatorname{Tr}}
\newcommand{\tr}{\operatorname{tr}}
\newcommand{\Rsq}{R_{\mu\nu\rho\sigma}R^{\mu\nu\rho\sigma}}
\newcommand{\Ricsq}{R_{\mu\nu}R^{\mu\nu}}
\newcommand{\Weylsq}{C_{\mu\nu\rho\sigma}C^{\mu\nu\rho\sigma}}
\newcommand{\dd}{\mathrm{d}}
\newcommand{\ii}{\mathrm{i}}
\newcommand{\erfi}{\operatorname{erfi}}

%% ---- annotation commands --------------------------------------------
\newcommand{\fromlit}[1]{\textcolor{blue!70!black}{\textsf{[#1]}}}
\newcommand{\oursign}{\textcolor{teal!80!black}{\textsf{[our convention]}}}
\newcommand{\gap}{\textcolor{red!80!black}{\textsf{[GAP]}}}
\newcommand{\needscheck}{\textcolor{orange!80!black}{\textsf{[needs check]}}}
\newcommand{\verified}[1]{%
  \par\smallskip\noindent
  \colorbox{green!10}{\parbox{\dimexpr\linewidth-2\fboxsep}{%
    \textcolor{green!50!black}{\textbf{[Verified]}} #1}}%
  \par\smallskip}
\newcommand{\signcorrection}[1]{%
  \par\noindent\colorbox{red!10}{%
    \parbox{\dimexpr\linewidth-2\fboxsep}{%
      \textcolor{red!50!black}{\textbf{[Convention note]}}
      \textcolor{red!50!black}{#1}%
    }%
  }\par%
}

%% ---- page style -----------------------------------------------------
\pagestyle{fancy}
\fancyhf{}
\fancyhead[L]{\small\textsc{MT-1 (v2): Literature for BH Entropy from SCT}}
\fancyhead[R]{\small\thepage}
\renewcommand{\headrulewidth}{0.4pt}

%% ---- hyperref -------------------------------------------------------
\hypersetup{
  colorlinks=true,
  linkcolor=blue!60!black,
  citecolor=green!40!black,
  urlcolor=blue!60!black,
  pdftitle={MT-1 v2: Literature Extraction for BH Entropy from SCT Spectral Action},
  pdfauthor={David Alfyorov},
}

%% ---- title ----------------------------------------------------------
\title{\textbf{MT-1: Literature Extraction \& Convention Fixing (v2)\\[4pt]
       for Black Hole Entropy from the Spectral Action}\\[8pt]
       \large Wald--Noether Charge, Sen Logarithmic Correction,\\
       Four Laws, and Gap Tracker}
\author{David Alfyorov}
\date{April 2026 --- Working Document (v2)}

%%======================================================================
\begin{document}
\maketitle

\begin{center}
\small\textit{Credits: Igor Shnyukov contributed research-assistance
and workflow support.}
\end{center}

\begin{abstract}
This is the corrected v2 of the MT-1 literature extraction document.
The v1 contained a critical sign error in the Sen (2012) logarithmic
correction formula (vector and graviton signs reversed, Weyl/Dirac
fermion conflation), caught by the L-R audit.  This v2 corrects all
errors and reorganises the material into nine sub-tasks (L-1 through
L-9) covering: the Wald--Noether charge formula with explicit binormal
normalisation (L-1); the Jacobson--Myers formula with the $8\pi$ vs
$4\pi$ ambiguity resolved (L-2); the Sen (2012) logarithmic correction
with explicit convention notes (L-3); Stelle gravity benchmarks (L-4);
NCG first-law violation (L-5); LQG comparison (L-6); Iyer--Wald first
law (L-7); Wall second law conditions (L-8); and a complete gap tracker
for all eight gaps G1--G8 (L-9).

\medskip\noindent
\textbf{Main results (extracted, not derived):}
\begin{itemize}[nosep]
\item $S = A/(4G) + 13/(120\pi) + (37/24)\ln(A/\ell_P^2) + \mathcal{O}(1)$
\item $\alpha_C = 13/120$ (parameter-free, from NT-1b Phase~3)
\item $c_{\log} = 37/24 \approx 1.542$ (microcanonical, from Sen 2012 + SM content)
\item $R^2$ sector contributes zero on Schwarzschild ($R = 0$, $R_{\mu\nu} = 0$)
\end{itemize}
\end{abstract}

\tableofcontents
\newpage


%%======================================================================
\section{Conventions Inherited from NT-1 through NT-4b}
\label{sec:conventions}
%%======================================================================

\begin{convention}[Metric and curvature]\label{conv:metric}
Lorentzian signature $(-,+,+,+)$.  Riemann tensor
$R^\rho{}_{\sigma\mu\nu} = \partial_\mu\Gamma^\rho_{\nu\sigma}
- \partial_\nu\Gamma^\rho_{\mu\sigma} + \cdots$,
Ricci tensor $R_{\mu\nu} = R^\rho{}_{\mu\rho\nu}$,
Weyl tensor via the standard decomposition.
$\kappa^2 = 8\pi G$.
\end{convention}

\begin{convention}[SCT spectral action coefficients]\label{conv:sct}
\fromlit{NT-1b Phase~3.}
\begin{center}
\begin{tabular}{lll}
\toprule
Symbol & Value & Origin \\
\midrule
$\alpha_C$ & $13/120$ & Total Weyl${}^2$ coefficient (parameter-free) \\
$\alpha_R(\xi)$ & $2(\xi - 1/6)^2$ & Total $R^2$ coefficient \\
$F_1(0)$ & $\alpha_C/(16\pi^2) = 13/(1920\pi^2)$ & Weyl form factor at $z=0$ \\
$F_2(0,\xi)$ & $\alpha_R(\xi)/(16\pi^2)$ & Ricci form factor at $z=0$ \\
$\varphi(x)$ & $e^{-x/4}\sqrt{\pi/x}\,\erfi(\sqrt{x}/2)$ & Master function (entire) \\
\bottomrule
\end{tabular}
\end{center}
\end{convention}

\begin{convention}[SM field content]\label{conv:SM}
\fromlit{CPR 0805.2909; NT-1b Phase~3.}
\begin{center}
\begin{tabular}{lccl}
\toprule
Field & Spin & Count & Convention \\
\midrule
Real scalars (Higgs doublet) & 0 & $N_s = 4$ & Real components \\
Weyl fermions & $1/2$ & $N_f = 45$ & Weyl count \\
Dirac fermions & $1/2$ & $N_D = N_f/2 = 22.5$ & Dirac count \\
Gauge bosons & 1 & $N_v = 12$ & $SU(3)\times SU(2)\times U(1)$ \\
\bottomrule
\end{tabular}
\end{center}
\signcorrection{The Sen (2012) formula uses \textbf{Dirac} fermion count
$N_D = 22.5$ with coefficient 7 per Dirac field, \emph{not} the Weyl
count $N_f = 45$.  The coefficient per Weyl fermion is $7/2$.  The v1
document conflated these conventions.}
\end{convention}

\begin{convention}[Schwarzschild horizon data]\label{conv:schw}
\fromlit{Wald (1984), Chapter~6; MTW (1973).}
\begin{center}
\begin{tabular}{ll}
\toprule
Quantity & Value \\
\midrule
Horizon radius & $r_H = 2GM$ \\
Surface gravity & $\kappa = 1/(4GM)$ \\
Area & $A = 4\pi r_H^2 = 16\pi G^2 M^2$ \\
Hawking temperature & $T_H = \kappa/(2\pi) = 1/(8\pi GM)$ \\
Bekenstein--Hawking entropy & $S_{\mathrm{BH}} = A/(4G) = 4\pi G M^2$ \\
Euler characteristic & $\chi(S^2) = 2$ \\
\bottomrule
\end{tabular}
\end{center}
Curvature invariants: $R = 0$, $R_{\mu\nu} = 0$ (vacuum),
$C_{\mu\nu\rho\sigma} = R_{\mu\nu\rho\sigma}$ (Ricci-flat),
$\Weylsq\big|_{r_H} = 3/(4G^4 M^4)$.
\end{convention}


%%======================================================================
\section{L-1: Wald--Noether Charge Formula}
\label{sec:L1}
%%======================================================================

\fromlit{Wald (1993), Phys.\ Rev.\ D~\textbf{48}, R3427,
arXiv:gr-qc/9307038.}

\subsection{General statement}

For any diffeomorphism-invariant Lagrangian
$\mathcal{L}(g_{\mu\nu}, R_{\mu\nu\rho\sigma},
\nabla_\alpha R_{\mu\nu\rho\sigma}, \ldots)$,
the black hole entropy is the Noether charge:
\begin{equation}\label{eq:wald-general}
  \boxed{
  S_{\mathrm{Wald}} = -2\pi \oint_{\mathcal{H}}
  \frac{\partial\mathcal{L}}{\partial R_{\mu\nu\rho\sigma}}\,
  \hat{\epsilon}_{\mu\nu}\,\hat{\epsilon}_{\rho\sigma}\,\dd A.}
\end{equation}

\subsection{Binormal normalisation convention}
\fromlit{Wald (1993), eqs.~(12)--(14); Iyer--Wald (1994), eq.~(8).}

The binormal $\hat{\epsilon}_{\mu\nu}$ to the bifurcation
2-surface $\mathcal{H}$ lies in the $(t,r)$ plane:
\begin{equation}\label{eq:binormal}
  \hat{\epsilon}_{\mu\nu} = 2\,n_{[\mu}\,s_{\nu]},
  \qquad
  \hat{\epsilon}_{\mu\nu}\hat{\epsilon}^{\mu\nu} = -2,
\end{equation}
where $n^\mu$ is the future-directed unit timelike normal
($n_\mu n^\mu = -1$) and $s^\mu$ is the outward unit spacelike
normal ($s_\mu s^\mu = +1$) to $\mathcal{H}$.

\begin{remark}[Sign of the normalisation]
The Lorentzian signature $(-,+,+,+)$ gives
$\hat{\epsilon}_{\mu\nu}\hat{\epsilon}^{\mu\nu} = -2$ (negative).
Some references using $(+,-,-,-)$ obtain $+2$.  The overall sign in
the Wald formula~\eqref{eq:wald-general} is adjusted accordingly:
the combination $(-2\pi)\times(-2) = +4\pi$ produces the positive
Bekenstein--Hawking entropy from the Einstein--Hilbert term.
\end{remark}

\subsection{Einstein--Hilbert contribution}
\fromlit{Wald (1993), eq.~(14).}

For $\mathcal{L}_{\mathrm{EH}} = R/(16\pi G)$:
\begin{equation}
  \frac{\partial\mathcal{L}_{\mathrm{EH}}}
  {\partial R_{\mu\nu\rho\sigma}}
  = \frac{1}{16\pi G}\,g^{\mu[\rho}\,g^{\sigma]\nu},
  \qquad
  g^{\mu[\rho}\,g^{\sigma]\nu}\,
  \hat{\epsilon}_{\mu\nu}\,\hat{\epsilon}_{\rho\sigma} = -2.
\end{equation}
Hence
\begin{equation}\label{eq:S-EH}
  S_{\mathrm{EH}} = -2\pi \cdot \frac{1}{16\pi G}
  \cdot (-2) \cdot A = \frac{A}{4G}.
\end{equation}

\subsection{Nonlocal Lagrangians}
\fromlit{Conroy--Edholm--Mazumdar (2017), arXiv:1503.05568, Section~5;
Buoninfante--Mazumdar (2020), arXiv:1903.01542, Section~3.}

When the Lagrangian contains $f(\Box)\,I$ where $I$ is a curvature
invariant, the Wald formula generalises via integration by parts.
For astrophysical BH with $r_H \gg 1/\Lambda$, the on-horizon
evaluation uses $f(\Box) \to f(0)$:
\begin{equation}
  \frac{\Box}{\Lambda^2}\bigg|_{\mathcal{H}} \sim
  \frac{\kappa_H^2}{\Lambda^2} \sim
  \frac{1}{(GM\Lambda)^2} \ll 1.
\end{equation}
Nonlocal corrections to the Wald entropy are suppressed by
$(M_{\mathrm{Pl}}/M)^2 \cdot (\Lambda/M_{\mathrm{Pl}})^2$.
\gap\ Full nonlocal Wald formula needed only for micro-BH (see G1).


%%======================================================================
\section{L-2: Jacobson--Myers Formula}
\label{sec:L2}
%%======================================================================

\fromlit{Jacobson--Myers (1993), Phys.\ Rev.\ Lett.~\textbf{70}, 3684;
Jacobson--Kang--Myers (1994), arXiv:gr-qc/9312023.}

\subsection{Statement}

For a Lagrangian $\mathcal{L} = \gamma\,\Rsq$ (Gauss--Bonnet-type
in the Riemann-squared sector), the Wald entropy correction on
a background with $R_{\mu\nu} = 0$ is
\begin{equation}\label{eq:JM-formula}
  \boxed{
  \delta S_{\mathrm{JM}} = 8\pi\gamma\,\chi(\mathcal{H}).}
\end{equation}

\subsection{The $8\pi$ vs $4\pi$ ambiguity}

\signcorrection{Some references quote $\delta S = 4\pi\gamma\,
\chi(\mathcal{H})$ without the factor of~2.  This arises from
different conventions for the Gauss--Bonnet coupling:
\begin{itemize}[nosep]
\item \textbf{Convention A} (Jacobson--Myers): $\mathcal{L} =
  \gamma\,\Rsq$.  The Wald variation gives a factor
  $2\gamma\,R_{\mu\alpha\nu\beta}\,\hat{\epsilon}^{\mu\nu}
  \hat{\epsilon}^{\alpha\beta}$, and the Gauss--Bonnet theorem
  converts the integral to $8\pi\gamma\,\chi(\mathcal{H})$.
\item \textbf{Convention B}: $\mathcal{L} = \tilde{\gamma}\,E_4$
  where $E_4 = \Rsq - 4\Ricsq + R^2$ is the full Euler density.
  On Ricci-flat backgrounds $E_4 = \Rsq$, so $\tilde{\gamma} =
  \gamma$, but the Wald variation of $E_4$ has different
  combinatorial factors.
\end{itemize}
For our application on \textbf{Ricci-flat} backgrounds:
$\Weylsq = \Rsq$, and the coefficient of $\Rsq$ in the SCT
action is $\gamma = \alpha_C/(16\pi^2)$.  Using Convention A:
\begin{equation}\label{eq:delta-S-C2}
  \delta S_{C^2} = 8\pi \cdot \frac{\alpha_C}{16\pi^2} \cdot 2
  = \frac{\alpha_C}{\pi}
  = \frac{13}{120\pi}
  \approx 0.03449.
\end{equation}
This is a \textbf{topological} correction, independent of the
horizon area~$A$.}

\subsection{$R^2$ contribution}
\fromlit{Jacobson--Myers (1993); Padmanabhan (2010), Chapter~15.}

For $\mathcal{L}_{R^2} = \beta\,R^2$ with
$\beta = \alpha_R(\xi)/(16\pi^2)$:
\begin{equation}
  \frac{\partial(R^2)}{\partial R_{\mu\nu\rho\sigma}}
  = 2R\,g^{\mu[\rho}\,g^{\sigma]\nu}.
\end{equation}
On Schwarzschild, $R = 0$, so $\delta S_{R^2}|_{\mathrm{Schw}} = 0$.
The $R^2$ sector contributes only on backgrounds with $R \neq 0$
(de~Sitter, Reissner--Nordstr\"om--de~Sitter, etc.).

\subsection{Weyl binormal contraction}
\fromlit{Padmanabhan (2010), eq.~(15.62).}

In the orthonormal frame adapted to the $(t,r)$ binormal:
\begin{equation}\label{eq:C-binormal}
  C_{\mu\nu\rho\sigma}\,\hat{\epsilon}^{\mu\nu}
  \hat{\epsilon}^{\rho\sigma}\big|_{\mathcal{H}}
  = 4\,C_{\hat{0}\hat{1}\hat{0}\hat{1}}\big|_{r_H}
  = 4 \cdot \Bigl(-\frac{2GM}{r_H^3}\Bigr)
  = -\frac{1}{G^2 M^2}.
\end{equation}
The uniformity on $S^2$ (spherical symmetry of the Schwarzschild
Weyl tensor) ensures the integrand in the Wald formula is constant
over $\mathcal{H}$, so $\oint_{\mathcal{H}} \to A \times (\cdots)$.


%%======================================================================
\section{L-3: Sen (2012) Logarithmic Correction Formula}
\label{sec:L3}
%%======================================================================

\fromlit{Sen (2012), arXiv:1205.0971, eq.~(1.2) and Table~1.}

\subsection{The formula}

For a 4D Schwarzschild BH in Einstein gravity coupled to $N_s$~real
scalars, $N_F$~massless \textbf{Dirac} fermions, and $N_V$~massless
vectors, the one-loop logarithmic correction to the microcanonical
entropy is:
\begin{equation}\label{eq:sen-formula}
  \boxed{
  C_{\mathrm{local}} = \frac{1}{90}\bigl[
    2\,N_s + 7\,N_D - 26\,N_V + 424
  \bigr],}
\end{equation}
and the logarithmic entropy correction coefficient is
\begin{equation}\label{eq:clog-def}
  \boxed{
  c_{\log} = \frac{C_{\mathrm{local}}}{2}
  = \frac{1}{180}\bigl[
    2\,N_s + 7\,N_D - 26\,N_V + 424
  \bigr].}
\end{equation}
The entropy acquires the correction
$\Delta S = c_{\log}\,\ln(A_H/\ell_P^2)$.

\signcorrection{%
\textbf{Convention note (v1 error corrected):}
\begin{enumerate}[nosep]
\item Sen's $C_{\mathrm{local}}$ uses $1/90$; our $c_{\log} =
  C_{\mathrm{local}}/2$ uses $1/180$.  This factor of~2 arises
  from the relationship between the heat kernel $a_4$ coefficient
  and the entropy logarithm.
\item $N_D$ counts \textbf{Dirac} fermions (not Weyl).
  The SM has $N_D = 22.5$ Dirac fermions ($= 45$ Weyl$/2$).
  The coefficient is $+7$ per Dirac, equivalently $+7/2$ per Weyl.
\item The vector coefficient is $-26$ per vector (\textbf{negative}),
  due to Faddeev--Popov ghost dominance in the Kretschner term of
  the $a_4$ Seeley--DeWitt coefficient.
\item The graviton contribution is $+424$ (\textbf{positive}),
  including zero-mode corrections.
\end{enumerate}
The v1 document had the vector and graviton signs reversed
($+26\,N_V - 424$ instead of $-26\,N_V + 424$) and used the Weyl
count with coefficient~7.  All three errors cancelled partially but
produced the wrong numerical result ($211/180$ instead of $37/24$).}

\subsection{Individual field contributions}

\begin{center}
\begin{tabular}{lcc}
\toprule
Field & Contribution to $180\,c_{\log}$ & Sign \\
\midrule
Real scalar & $+2$ per field & positive \\
Dirac fermion & $+7$ per field & positive \\
Massless vector & $-26$ per field & \textbf{negative} \\
Graviton & $+424$ (total) & \textbf{positive} \\
\bottomrule
\end{tabular}
\end{center}

\subsection{Pure gravity cross-check}
\fromlit{Sen (2012), Section~3.}

Setting $N_s = N_D = N_V = 0$:
\begin{equation}
  c_{\log}^{\mathrm{grav}} = \frac{424}{180}
  = \frac{106}{45} \approx 2.356
  \quad (\text{positive}).
\end{equation}
Sen explicitly states pure gravity gives a positive logarithmic
correction.  The v1 formula gave $-424/180$ (negative), which is
the definitive sign of the error.

\subsection{SCT prediction for $c_{\log}$}

Using the SM content from Convention~\ref{conv:SM}
($N_s = 4$, $N_D = 22.5$, $N_V = 12$):
\begin{equation}\label{eq:clog-sct}
  \boxed{
  c_{\log} = \frac{1}{180}\bigl[
    2 \cdot 4 + 7 \cdot 22.5 - 26 \cdot 12 + 424
  \bigr]
  = \frac{277.5}{180}
  = \frac{37}{24}
  \approx 1.5417.}
\end{equation}

\begin{proof}[Arithmetic verification]
$2 \times 4 = 8$;\quad $7 \times 22.5 = 157.5$;\quad
$26 \times 12 = 312$;\quad graviton: $424$.\\
Numerator: $8 + 157.5 - 312 + 424 = 277.5$.\\
$277.5/180 = 277.5/180$.
Simplify: $277.5 = 1110/4$, so $c_{\log} = 1110/(4 \times 180)
= 1110/720 = 37/24$.
\end{proof}

\begin{remark}[Matter vs.\ graviton decomposition]
Graviton: $+424/180 = 106/45 \approx +2.356$ (positive: gravity
\emph{increases} the correction).
Matter: $(8 + 157.5 - 312)/180 = -146.5/180 \approx -0.814$
(net negative: vectors dominate).
Total: $c_{\log} = 37/24 > 0$.
\end{remark}

\subsection{Weyl fermion convention (alternative form)}

Substituting $N_D = N_W/2$ where $N_W$ is the Weyl fermion count:
\begin{equation}\label{eq:sen-weyl}
  c_{\log} = \frac{1}{180}\bigl[
    2\,N_s + \tfrac{7}{2}\,N_W - 26\,N_V + 424
  \bigr].
\end{equation}
Check: $N_W = 45$, coefficient $= 7/2 \times 45 = 157.5$. Consistent.

\subsection{Independent confirmation}
\fromlit{arXiv:2104.14902, eq.~(5).}

The coefficient $\gamma$ in the heat kernel is:
\begin{equation}
  \gamma = \frac{1}{11520\pi^2}\bigl[
    2\,N_s + 7\,N_D - 26\,N_V + 424
  \bigr],
\end{equation}
and $c_{\log} = 64\pi^2\,\gamma$, which reproduces
eq.~\eqref{eq:clog-def}.

\subsection{Microcanonical vs singlet convention}

The value $c_{\log} = 37/24$ is the \textbf{microcanonical}
coefficient (no zero-mode subtraction).  The rotationally invariant
(\textbf{singlet}) entropy subtracts $N_R = 3$ rotation generators:
\begin{equation}
  c_{\log}^{\mathrm{singlet}}
  = \frac{C_{\mathrm{local}} - 3}{2}
  = \frac{277.5/90 - 3}{2}
  = \frac{37/24 - 3/2}{1}
  = \frac{37 - 36}{24}
  = \frac{1}{24}
  \approx 0.0417.
\end{equation}
We use the microcanonical convention throughout.


%%======================================================================
\section{L-4: Stelle Gravity Wald Entropy on Schwarzschild}
\label{sec:L4}
%%======================================================================

\fromlit{Stelle (1977), Phys.\ Rev.\ D~\textbf{16}, 953;
Gen.\ Rel.\ Grav.~\textbf{9} (1978), 353.}

\subsection{Stelle action}

The Stelle action in 4D is
\begin{equation}\label{eq:stelle}
  S_{\mathrm{Stelle}} = \int\dd^4 x\,\sqrt{-g}\,
  \biggl[\frac{R}{16\pi G}
  + \alpha_{\mathrm{St}}\,\Weylsq
  + \beta_{\mathrm{St}}\,R^2\biggr],
\end{equation}
with $\alpha_{\mathrm{St}}$ and $\beta_{\mathrm{St}}$ as free
couplings.

\subsection{Wald entropy}

On Schwarzschild ($R = 0$, $R_{\mu\nu} = 0$):
\begin{equation}
  S_{\mathrm{Wald}}^{\mathrm{Stelle}}
  = \frac{A}{4G} + \frac{\alpha_{\mathrm{St}}}{\pi}\,
  \chi(\mathcal{H})
  = \frac{A}{4G} + \frac{2\alpha_{\mathrm{St}}}{\pi}.
\end{equation}
The $R^2$ term contributes zero.  This is the local HD gravity
benchmark.

\subsection{SCT as a specific Stelle limit}

In the local limit ($\Box/\Lambda^2 \to 0$), the SCT spectral
action reduces to Stelle gravity with
$\alpha_{\mathrm{St}} = \alpha_C/(16\pi^2) = 13/(1920\pi^2)$,
$\beta_{\mathrm{St}} = \alpha_R(\xi)/(16\pi^2)$.  The Wald
entropy then gives:
\begin{equation}
  \delta S_{C^2}^{\mathrm{SCT}} = \frac{2\alpha_C}{16\pi^2 \cdot \pi}
  \cdot \chi = \frac{\alpha_C}{\pi} = \frac{13}{120\pi},
\end{equation}
reproducing eq.~\eqref{eq:delta-S-C2}. This serves as a
consistency check.


%%======================================================================
\section{L-5: NCG First Law Violation}
\label{sec:L5}
%%======================================================================

\fromlit{arXiv:2509.23423 (2025).}

\subsection{The claim}

A 2025 analysis within noncommutative geometry (NCG) shows that the
conventional first law $\dd M = (\kappa/2\pi)\,\dd S + \cdots$ is
violated when the entropy is identified with the Bekenstein--Hawking
area formula $S = A/(4G)$.  The violation arises because the NCG
action contains higher-curvature terms that shift the entropy away
from the area law.

\subsection{Relevance to SCT}

In SCT, the first law is \textbf{maintained} precisely because we use
the \textbf{Wald entropy} (Noether charge), not the area entropy.
The Iyer--Wald theorem (Section~\ref{sec:L7}) guarantees the first
law for any diffeomorphism-invariant Lagrangian when the correct
(Wald) entropy is used.  The NCG ``violation'' is an artefact of
using $A/(4G)$ instead of $S_{\mathrm{Wald}}$.

\begin{remark}[Lesson for SCT]
The NCG result reinforces that in any HD gravity theory (including
SCT), one \textbf{must} use the Wald entropy to preserve the first
law.  The area entropy is not the correct thermodynamic potential.
\end{remark}


%%======================================================================
\section{L-6: LQG Logarithmic Correction Comparison}
\label{sec:L6}
%%======================================================================

\fromlit{Kaul--Majumdar (2000), Phys.\ Rev.\ Lett.~\textbf{84}, 5255,
arXiv:gr-qc/0002040;
Engle--Noui--Perez (2010), Phys.\ Rev.\ Lett.~\textbf{105}, 031302,
arXiv:0905.3168.}

\subsection{LQG result}

In loop quantum gravity (isolated horizon framework), the
logarithmic correction to BH entropy is:
\begin{equation}\label{eq:clog-lqg}
  c_{\log}^{\mathrm{LQG}} = -\frac{3}{2}.
\end{equation}
This result arises from the $SU(2)$ Chern--Simons state counting
on the horizon, where the $-3/2$ is related to the zero mode of
the $SU(2)$ gauge field (specifically, $3$ rotational modes
subtracted with a $-1/2$ factor from the logarithm of the
partition function).

\subsection{Comparison table}

\begin{center}
\begin{tabular}{lccl}
\toprule
Approach & $c_{\log}$ & Sign & Reference \\
\midrule
\textbf{SCT (this work)} & $\mathbf{37/24 \approx +1.542}$
  & \textbf{positive} & parameter-free \\
LQG (isolated horizons) & $-3/2 = -1.500$
  & negative & Kaul--Majumdar (2000) \\
LQG (Engle--Noui--Perez) & $-3/2$
  & negative & arXiv:0905.3168 \\
Pure gravity (Euclidean) & $+106/45 \approx +2.356$
  & positive & Sen (2012) \\
String theory (extremal) & model-dependent
  & varies & Sen (2012) \\
IDG (infinite derivatives) & same as Sen formula
  & positive & Modesto (2012) \\
Asymptotic safety & not well-defined
  & --- & (running $G$) \\
\bottomrule
\end{tabular}
\end{center}

\begin{remark}[Discriminating observable]
The SCT prediction $c_{\log} = +37/24$ has the \textbf{opposite sign}
from the LQG prediction $c_{\log}^{\mathrm{LQG}} = -3/2$.
Sen (2012) notes this disagreement explicitly, observing that the
Euclidean gravity computation (which SCT inherits) gives a positive
$c_{\log}$, while LQG gives a negative one.  This sign difference
is a potential observational discriminant between the two frameworks,
though measuring logarithmic corrections to BH entropy is far beyond
current experimental capability.
\end{remark}


%%======================================================================
\section{L-7: Iyer--Wald First Law}
\label{sec:L7}
%%======================================================================

\fromlit{Iyer--Wald (1994), Phys.\ Rev.\ D~\textbf{50}, 846,
arXiv:gr-qc/9403028.}

\subsection{Statement}

\begin{theorem}[Iyer--Wald, 1994]\label{thm:iyer-wald}
For any diffeomorphism-invariant Lagrangian
$\mathbf{L} = L\,\boldsymbol{\epsilon}$, where
$\boldsymbol{\epsilon}$ is the volume form, and for any
one-parameter family of stationary, axisymmetric black hole
solutions with bifurcate Killing horizon, the first law holds:
\begin{equation}\label{eq:first-law}
  \boxed{
  \delta M = \frac{\kappa}{2\pi}\,\delta S_{\mathrm{Wald}}
  + \Omega_H\,\delta J + \Phi_H\,\delta Q,}
\end{equation}
where $S_{\mathrm{Wald}}$ is the Noether charge entropy~\eqref{eq:wald-general}, $\kappa$ is the surface gravity, $\Omega_H$
the angular velocity of the horizon, $\Phi_H$ the electric
potential, and $M$, $J$, $Q$ are the ADM mass, angular momentum,
and charge.
\end{theorem}

\subsection{Derivation outline}
\fromlit{Iyer--Wald (1994), Sections~III--IV.}

The proof uses the Noether current $\mathbf{J}[\xi]$ and Noether
charge $\mathbf{Q}[\xi]$ associated with the Killing field
$\xi^a = t^a + \Omega_H\,\phi^a$:
\begin{enumerate}[nosep]
\item Construct the symplectic current
  $\boldsymbol{\omega}(\delta_1 g, \delta_2 g)$ from the
  Lagrangian variation.
\item Apply Stokes' theorem between spatial infinity (where
  $\delta M$, $\delta J$, $\delta Q$ are read off) and the
  bifurcation surface $\mathcal{H}$.
\item On $\mathcal{H}$, the Killing field $\xi^a$ vanishes, and
  the Noether charge integral reduces to the Wald
  entropy~\eqref{eq:wald-general}.
\end{enumerate}
No energy conditions are required.  The proof is purely geometric.

\subsection{Status in SCT}

The SCT spectral action is diffeomorphism-invariant (it is
constructed from the spectral triple, which is a geometric object).
Therefore the Iyer--Wald first law applies automatically to the
Wald entropy of any stationary black hole solution in SCT.
No additional verification is needed beyond confirming
diffeomorphism invariance.


%%======================================================================
\section{L-8: Wall (2015) Second Law for HD Gravity}
\label{sec:L8}
%%======================================================================

\fromlit{Wall (2015), Int.\ J.\ Mod.\ Phys.\ D~\textbf{24}, 1544014,
arXiv:1504.08040;
Wall (2024), arXiv:2402.05411.}

\subsection{Statement}

\begin{theorem}[Wall, 2015]\label{thm:wall}
For a diffeomorphism-invariant theory of gravity in $d \geq 4$
dimensions, the generalised second law
\begin{equation}
  \delta S_{\mathrm{gen}} = \delta S_{\mathrm{Wald}}
  + \delta S_{\mathrm{outside}} \geq 0
\end{equation}
holds provided three conditions are satisfied:
\begin{enumerate}[nosep]
\item \textbf{Well-posed initial value problem} for the
  gravitational field equations.
\item \textbf{Null energy condition (NEC)} for the matter sector:
  $T_{\mu\nu}\,k^\mu k^\nu \geq 0$ for all null $k^\mu$.
\item \textbf{Linearised stability:} the entropy functional
  satisfies certain positivity conditions under linearised
  perturbations of the horizon geometry.
\end{enumerate}
\end{theorem}

\subsection{Status in SCT}

\begin{enumerate}[nosep]
\item \textbf{Well-posedness:} The SCT field equations contain
  nonlocal (infinite-derivative) terms.  The entire-function
  property of $F_1(z)$ (NT-2) prevents additional propagating
  degrees of freedom beyond those of Stelle gravity.  However,
  formal well-posedness for the full nonlocal equations in the
  sense required by Wall is an open problem.
  \gap\ (Connects to MR-4.)

\item \textbf{NEC:} The SM matter sector satisfies the NEC
  classically.  Quantum violations (Casimir-type) are
  $\mathcal{O}(\hbar)$ and do not affect the semiclassical
  second law.

\item \textbf{Linearised stability:} The spin-2 propagator
  $\Pi_{\mathrm{TT}}(z)$ has ghost poles (MR-2 ghost catalogue).
  These complex poles may violate Wall's stability condition.
  The Lee--Wick/fakeon prescription (FK analysis) may render
  the ghosts harmless in the macroscopic regime, but this is
  not rigorously established for the second law.
  \gap\ (Connects to MR-2.)
\end{enumerate}

\begin{tcolorbox}[colback=orange!5, colframe=orange!60!black]
\textbf{Verdict:} The second law in SCT is \textbf{CONDITIONAL}
on the ghost resolution (MR-2) and well-posedness (MR-4).
For macroscopic BH where ghost effects are suppressed by
$(M_{\mathrm{Pl}}/\Lambda)^2$, the second law is expected to
hold, but a rigorous proof is lacking.
\end{tcolorbox}


%%======================================================================
\section{L-9: Gap Tracker}
\label{sec:L9}
%%======================================================================

All eight gaps identified during the literature extraction are tracked
below.  Status labels: \textcolor{green!50!black}{\textbf{RESOLVED}},
\textcolor{orange!80!black}{\textbf{OPEN}},
\textcolor{red!80!black}{\textbf{BLOCKING}},
\textcolor{blue!60!black}{\textbf{DEFERRED}}.

\begin{enumerate}[label=\textbf{G\arabic*:},leftmargin=3em]

\item \textbf{Nonlocal Wald formula.}
  The full Wald formula for a Lagrangian $\int F(\Box)\,C^2$ requires
  integration by parts of the form factor away from the local limit.
  For astrophysical BH ($r_H \gg 1/\Lambda$), the local limit
  $F(\Box) \to F(0)$ suffices and the standard Wald formula applies.
  For micro-BH ($r_H \sim 1/\Lambda$), the full nonlocal formula is
  essential.\\
  \textbf{Status:} \textcolor{orange!80!black}{\textbf{OPEN}}
  (non-blocking for astrophysical BH).\\
  \textbf{Resolution path:} Use the Conroy--Edholm--Mazumdar (2017)
  framework for nonlocal Wald entropy.  The entire-function property
  of $F_1(z)$ from NT-2 ensures no spurious divergences.

\item \textbf{Weyl binormal contraction in Kruskal coordinates.}
  The value $C_{\mu\nu\rho\sigma}\,\hat{\epsilon}^{\mu\nu}
  \hat{\epsilon}^{\rho\sigma}|_{\mathcal{H}} = -1/(G^2 M^2)$
  is computed in the orthonormal frame.  A cross-check in
  regular Kruskal coordinates would strengthen confidence.\\
  \textbf{Status:} \textcolor{orange!80!black}{\textbf{OPEN}}
  (non-blocking; orthonormal frame result is standard).\\
  \textbf{Resolution path:} Direct computation in Kruskal
  coordinates, verifying $C_{UVUV}$ on the bifurcation surface.

\item \textbf{Second law and ghost sector.}
  Wall's second law requires well-posedness and linearised stability.
  The SCT ghost poles (MR-2: 2~real + 3~Lee--Wick pairs) may violate
  these conditions.\\
  \textbf{Status:} \textcolor{red!80!black}{\textbf{BLOCKING}}
  for the second law.  Does not affect the Wald entropy formula
  or $c_{\log}$.\\
  \textbf{Resolution path:} CONDITIONAL on MR-2 ghost resolution.
  If fakeon prescription is adopted, the macroscopic regime
  ($M \gg M_{\mathrm{Pl}}$) should satisfy Wall's conditions.

\item \textbf{Entanglement entropy route.}
  The spectral triple approach to BH entropy via entanglement
  across the horizon has not been computed explicitly in SCT.
  This is an alternative derivation route, not a prerequisite
  for the Wald entropy result.\\
  \textbf{Status:} \textcolor{blue!60!black}{\textbf{DEFERRED}}
  to future work.\\
  \textbf{Resolution path:} Requires near-horizon spectral
  geometry of the SCT Dirac operator.

\item \textbf{Kerr generalisation.}
  Kerr is vacuum ($R = 0$, $R_{\mu\nu} = 0$), so the $R^2$
  sector still vanishes and $\chi(S^2) = 2$ still applies.
  The topological correction $\alpha_C/\pi$ is unchanged.
  Frame-dragging modifies the Weyl binormal contraction, but
  the \emph{integrated} Wald entropy depends only on
  $\chi(\mathcal{H})$ via the Gauss--Bonnet theorem.\\
  \textbf{Status:} \textcolor{green!50!black}{\textbf{RESOLVED}}
  for the Wald entropy.  The $c_{\log}$ for Kerr requires
  separate analysis (Sen 2012 addresses this).\\
  \textbf{Resolution path:} None needed for Schwarzschild.
  Kerr $c_{\log}$ deferred.

\item \textbf{Sen formula: graviton zero modes.}
  The $+424$ graviton contribution in Sen's formula includes
  subtleties with zero modes on the Euclidean Schwarzschild
  instanton.  In SCT, the graviton propagator is modified by
  $\Pi_{\mathrm{TT}}(z)$, which shifts the pole structure.
  This may alter the zero-mode counting if the spectral
  cutoff $\Lambda$ is close to the Planck scale.\\
  \textbf{Status:} \textcolor{orange!80!black}{\textbf{OPEN}}
  (non-blocking for $\Lambda \ll M_{\mathrm{Pl}}$).\\
  \textbf{Resolution path:} Check whether the modified
  propagator changes the zero-mode contribution.  Expected
  to be negligible for $\Lambda \ll M_{\mathrm{Pl}}$.

\item \textbf{Modified evaporation endpoint.}
  Late-stage BH evaporation ($T_H \sim \Lambda$ at
  $M \sim M_{\mathrm{Pl}}^2/\Lambda$) requires the full nonlocal
  Regge--Wheeler equation and connects to MR-9 (singularity
  resolution).\\
  \textbf{Status:} \textcolor{blue!60!black}{\textbf{DEFERRED}}
  to MR-9.\\
  \textbf{Resolution path:} Solve the modified RW equation
  with full $\Pi_{\mathrm{TT}}(\omega^2/\Lambda^2)$ factor.

\item \textbf{Information paradox.}
  The entire-function property of the spectral action may provide
  a natural resolution through modified near-horizon physics
  (no trans-Planckian modes).\\
  \textbf{Status:} \textcolor{blue!60!black}{\textbf{DEFERRED}}
  to MR-9.\\
  \textbf{Resolution path:} Requires completion of the
  evaporation endpoint analysis (G7) and singularity resolution.

\end{enumerate}

\begin{center}
\begin{tabular}{clll}
\toprule
\textbf{Gap} & \textbf{Description} & \textbf{Status} & \textbf{Blocks} \\
\midrule
G1 & Nonlocal Wald formula & OPEN & Micro-BH only \\
G2 & Weyl binormal in Kruskal & OPEN & Nothing \\
G3 & Second law + ghosts & \textbf{BLOCKING} & Second law \\
G4 & Entanglement entropy & DEFERRED & Alternative route \\
G5 & Kerr generalisation & RESOLVED & --- \\
G6 & Graviton zero modes & OPEN & Micro-BH regime \\
G7 & Evaporation endpoint & DEFERRED (MR-9) & Late-stage physics \\
G8 & Information paradox & DEFERRED (MR-9) & Foundational \\
\bottomrule
\end{tabular}
\end{center}


%%======================================================================
\section{Full BH Entropy Formula in SCT (Literature Assembly)}
\label{sec:full-formula}
%%======================================================================

Combining the Wald entropy (L-1, L-2) with the logarithmic
correction (L-3):

\begin{tcolorbox}[colback=blue!5!white, colframe=blue!50!black,
  title={\textbf{MT-1 Main Formula (to be verified in derivation phase)}}]
\begin{equation}\label{eq:full-entropy}
  S = \underbrace{\frac{A}{4G}}_{\text{Bekenstein--Hawking}}
  + \underbrace{\frac{13}{120\pi}}_{\text{Wald ($C^2$, topological)}}
  + \underbrace{\frac{37}{24}\,
    \ln\!\left(\frac{A}{\ell_P^2}\right)}_{\text{one-loop (Sen)}}
  + \mathcal{O}(1).
\end{equation}
\end{tcolorbox}

\noindent
Properties of each term:
\begin{enumerate}[nosep]
\item \textbf{$A/(4G)$}: classical area law from the Einstein--Hilbert
  sector (Wald entropy).  Dominant for all $A \gg \ell_P^2$.
\item \textbf{$13/(120\pi) \approx 0.035$}: tree-level correction
  from the Weyl-squared sector; topological (depends on
  $\chi(\mathcal{H}) = 2$, not on $A$); parameter-free via
  $\alpha_C = 13/120$.
\item \textbf{$(37/24)\ln(A/\ell_P^2) \approx 1.542\,
  \ln(A/\ell_P^2)$}: one-loop quantum correction from all SM
  fields + graviton; parameter-free via SM field content.
  Microcanonical convention.
\end{enumerate}

\medskip\noindent
\textbf{Hierarchy for $M = 10\,M_\odot$:}
$A/(4G) \approx 1.05 \times 10^{79}$;\quad
$\delta S_{\mathrm{Wald}} \approx 0.035$;\quad
$\delta S_{\log} \approx (37/24) \times \ln(16\pi G^2 M^2/\ell_P^2)
\approx 283$.\\
Ordering: $S_{\mathrm{BH}} \gg \delta S_{\log} \gg
\delta S_{\mathrm{Wald}}$.


%%======================================================================
\section{Master Reference Table}
\label{sec:refs}
%%======================================================================

\begin{center}
\begin{longtable}{p{0.7cm}p{5.5cm}p{3.5cm}p{3.8cm}}
\toprule
\textbf{\#} & \textbf{Citation} & \textbf{Key result}
  & \textbf{Use in MT-1} \\
\midrule
\endhead
R1 & Wald (1993), PRD~\textbf{48}, R3427; gr-qc/9307038
  & BH entropy = Noether charge
  & L-1: Wald formula \\
R2 & Iyer--Wald (1994), PRD~\textbf{50}, 846; gr-qc/9403028
  & First law for diff-inv.\ theories
  & L-7: First law \\
R3 & Jacobson--Myers (1993), PRL~\textbf{70}, 3684
  & Entropy of Lovelock gravity
  & L-2: JM formula \\
R4 & Jacobson--Kang--Myers (1994), gr-qc/9312023
  & $C^2$ and $R^2$ Wald entropy
  & L-2: Benchmark \\
R5 & Sen (2012), arXiv:1205.0971
  & $c_{\log}$ for non-extremal BH
  & L-3: $c_{\log} = 37/24$ \\
R6 & Wall (2015), IJMPD~\textbf{24}, 1544014; 1504.08040
  & Second law for HD gravity
  & L-8: Conditions \\
R7 & Wall (2024), arXiv:2402.05411
  & Updated second law
  & L-8: Supplement \\
R8 & Bardeen--Carter--Hawking (1973), CMP~\textbf{31}, 161
  & Four laws of BH mechanics
  & Zeroth, third laws \\
R9 & Hawking (1975), CMP~\textbf{43}, 199
  & Hawking radiation
  & Temperature \\
R10 & Stelle (1977), PRD~\textbf{16}, 953
  & Higher-derivative gravity
  & L-4: Benchmark \\
R11 & Buoninfante--Mazumdar (2020), arXiv:1903.01542
  & Nonlocal gravity + BH thermo
  & L-1: Local limit \\
R12 & Conroy--Edholm--Mazumdar (2017), arXiv:1503.05568
  & Nonlocal Wald entropy
  & L-1: Nonlocal corrections \\
R13 & arXiv:2509.23423 (2025)
  & NCG first law violation
  & L-5: Lesson for SCT \\
R14 & Kaul--Majumdar (2000), PRL~\textbf{84}, 5255
  & LQG $c_{\log} = -3/2$
  & L-6: Comparison \\
R15 & Engle--Noui--Perez (2010), PRL~\textbf{105}, 031302
  & LQG $c_{\log}$ confirmation
  & L-6: Comparison \\
R16 & arXiv:2104.14902
  & Heat kernel $\gamma$ coefficient
  & L-3: Cross-check \\
R17 & Christensen--Duff (1978), NPB~\textbf{154}, 301
  & Conformal anomaly coefficients
  & Supplementary \\
R18 & Padmanabhan (2010), \textit{Gravitation}
  & Weyl tensor on horizons
  & L-2: Binormal \\
R19 & Bombelli et al.\ (1986), PRD~\textbf{34}, 373
  & Entanglement entropy = area law
  & G4: Future route \\
R20 & Mann--Solodukhin (1998), hep-th/9709064
  & Conical singularity method
  & Alternative derivation \\
R21 & Connes--Chamseddine (2010), arXiv:1004.0464
  & Spectral action and gravity
  & Spectral framework \\
\bottomrule
\end{longtable}
\end{center}


%%======================================================================
\section{Notation Summary}
\label{sec:notation}
%%======================================================================

\begin{center}
\begin{tabular}{lll}
\toprule
Symbol & Definition & Value/Source \\
\midrule
$\alpha_C$ & Weyl${}^2$ coefficient & $13/120$ (NT-1b Phase~3) \\
$\alpha_R(\xi)$ & $R^2$ coefficient & $2(\xi - 1/6)^2$ \\
$c_{\log}$ & Log entropy correction (microcan.) & $37/24 \approx 1.542$ \\
$C_{\mathrm{local}}$ & Sen's heat kernel coefficient & $277.5/90 = 37/12$ \\
$\kappa$ & Surface gravity & $1/(4GM)$ (Schwarzschild) \\
$T_H$ & Hawking temperature & $1/(8\pi GM)$ \\
$r_H$ & Horizon radius & $2GM$ \\
$A$ & Horizon area & $16\pi G^2 M^2$ \\
$\hat{\epsilon}_{\mu\nu}$ & Bifurcation surface binormal &
  $\hat{\epsilon}_{ab}\hat{\epsilon}^{ab} = -2$ \\
$\chi(\mathcal{H})$ & Euler characteristic & $2$ for $S^2$ \\
$N_s$ & Real scalars (SM) & $4$ \\
$N_D$ & Dirac fermions (SM) & $22.5$ \\
$N_V$ & Vectors (SM) & $12$ \\
$N_f$ & Weyl fermions (SM) & $45 = 2\,N_D$ \\
\bottomrule
\end{tabular}
\end{center}


\vfill
\begin{center}
\small\textit{End of MT-1 Literature Extraction v2.\\
All formulas extracted from cited sources; no original derivations.\\
Sen formula corrected from v1 (L-R audit: vector/graviton signs,
Dirac convention).}
\end{center}

\end{document}
